\begin{abstract}
Industrial inspection, driving and quality control decide by rules: a scene is level one if \emph{any} level-one hazard is present, a car must stop for a red light \emph{or} an obstacle. Vision-language models are adapted to such tasks by fine-tuning on verdicts. We show that this silently erases facts. Once one fact decides the verdict, every other fact can take any value without changing the label; it receives no supervision and, more fundamentally, no information. We call these facts \emph{overexposed}: like the highlights of an overexposed photograph, they are lost while the picture still looks right, and the damage appears only when the exposure changes---when the rule is edited. We make this precise: the information a verdict carries about a fact equals the fact's \emph{exposure}, the probability that flipping it changes the verdict, times its entropy, and the error after a rule edit is bounded by each fact's exposure under the \emph{new} rule. We propose \lux, which reads atomic facts in one parallel pass, composes the verdict with an editable rule table, and is trained with an objective derived from this analysis. On three datasets, pre-registered interventions confirm the predicted overexposure: verdict training raises false positives on clipped facts by 1.7--4.6$\times$ while leaving visible ones unchanged. Against blind human labels of real regulation edits on industrial inspection, \lux is correct on 89\% of the affected scenes; a fine-tuned model told the new rule is correct on 26\%.
\end{abstract}
